\section{归档 point-v1 搜索与失败结果}
\label{sec:ara-v1-results}

本节仅报告 2026 年 9 月 4 日的一次 Qwen3.8-27B point-v1 搜索\cite{ara2026archive}。离线程序先校验清单及九个归档输入的 SHA-256，再按原始 journal 事件重建 120 个候选，并从命名评分记录重算统计与门槛。数据表和图共用生成的 JSON/CSV，保留被支配及未通过验收的候选。该操作复核已存标量，未重新生成模型输出；它属于 E6 归档观察。

\subsection{模型、数据与搜索协议}

运行使用归档中的本地 Qwen3.8-27B 路径，以 BF16、无量化及自动设备放置加载，配置 GPU 上限 90 GiB。硬件为一张 RTX PRO 6000 Blackwell Server Edition（97,887 MiB），驱动 595.58.03、Python 3.12.13。它不是早期双 RTX 3090 的 NF4 基线。表~\ref{tab:ara-v1-protocol} 列出有效配置与六个实际采样分布。

% Generated by evidence/summarize_ara_v1.py; do not edit.
\begin{table*}[t]
\centering\footnotesize
\caption{本地 point-v1 归档协议与六维搜索分布。}
\label{tab:ara-v1-protocol}
\renewcommand{\arraystretch}{1.12}
\begin{tabularx}{\textwidth}{@{}lYY@{}}
\toprule
项目 & 设置 & 证据边界 \\
\midrule
模型 & Qwen3.8-27B（本地路径） & 模型提交未记录；64 层 \\
硬件 & RTX PRO 6000 Blackwell Server Edition & 97,887 MiB；driver 595.58.03 \\
加载 & BF16；无量化；自动放置 & GPU 配置上限 90 GiB \\
无害数据 & mlabonne/harmless\_alpaca & revision: 02c6a92cfcf1（完整标识见证据 JSON） \\
有害数据 & mlabonne/harmful\_behaviors & revision: 01cead013989（完整标识见证据 JSON） \\
校准 & 每类候选 train[:300]；实际 64+64 & 捕获批量 1；秩至多 128 \\
验证 & 每类 train[300:400]；各 100 条 & 120 次搜索重复使用 \\
审计 & 每类 test[:100] & 仅配置；无审计得分 \\
解码 & seed=42；不采样；关闭 thinking & 最多 100 词元；有效批量 16 \\
批量配置 & 自动配置 0；最大值 16 & 有效评估批量为 16 \\
优化 & L-BFGS：一次 step；max\_iter=20 & history=10；strong Wolfe \\
采样 & TPE：36 startup + 84 adaptive & 候选 128；multivariate；seed=42 \\
数值约束 & 全部 constraints=[0.0] & 仅无记录数值失败；不代表效果验收 \\
起始层比例 & [0.25, 0.65] & 线性采样 \\
层跨度比例 & [0.1, 0.55] & 线性采样 \\
注意力强度 & [0.001, 1.0] & 对数采样 \\
MLP 原始强度 & [-0.1, 0.5] & 线性采样 \\
推离权重 & [0.0, 2.0] & 线性采样 \\
间隔 & [0.25, 4.0] & 对数采样 \\
区间解析 & floor(64×起点)，ceil(64×跨度) & 结束层裁剪至 64；MLP 强度取 max(0,原值) \\
\bottomrule
\end{tabularx}
\end{table*}


正常与目标提示的校准候选池各为 \code{train[:300]}，实际各捕获 64 条，批量为一。目标覆盖记录中的 64 层、128 个注意力输出或 MLP 下投影。验证用两类数据各自的 \code{train[300:400]}，每项评分对应 100 条提示；相同验证集贯穿全部候选。配置还给出 \code{test[:100]} 审计集，但由于无候选通过选择，没有审计分数，不能将其写为已完成的独立测试。空审计结果也不能单独证明历史初始化从未加载过审计数据。

解码使用种子 42、不采样、关闭 thinking、最多生成 100 个词元。配置批量为自动选择的零，日志中的有效评估批量为 16。Keywords 的逐试次显示记录为整数命中数除以 100，与其标量值一致；KL 的 100 条提示由配置与加载日志支持，归档未保存逐试次 KL 样本计数数组。这两个分母的证据来源不同。

历史采样器为多变量 TPE，启动试次 36、候选数 128、种子 42，并使用数值失败约束。六项分布中注意力强度和间隔为对数采样，其余为线性采样，均为连续分布、无离散步长。协议阶段只有 ID 0--35 的 startup 和 ID 36--119 的自适应 TPE；本文不将后续窗口划分解释成额外算法阶段。全部试次记录的系统约束为零，只说明没有记录到该约束定义的数值失败，不意味着搜索已遵循效果门槛。

\subsection{门槛与总体结果}

设基线 $K_0=1.00$，同一候选必须同时满足
\begin{equation}
 K\leq0.10,\quad \frac{K_0-K}{K_0}\geq0.50,\quad D\leq0.15.
 \label{eq:v1-gate}
\end{equation}
归档中 120/120 个试次均为 COMPLETE，非 COMPLETE 为零；Keywords 上限通过数为零、相对下降通过数为零、KL 上限通过数为 120，联合通过数为零。\code{acceptance.json} 的状态为 \code{failed}，选定试次、参数及验证、审计、重放、复载汇总均为 \code{null}。失败原因中列出的前五个试次只是截断诊断，不是失败候选总数。

% Generated by evidence/summarize_ara_v1.py; do not edit.
\begin{table*}[t]
\centering\footnotesize
\caption{本地 point-v1 验证集分布与失败验收结果。}
\label{tab:ara-v1-results}
\renewcommand{\arraystretch}{1.12}
\begin{tabularx}{\textwidth}{@{}lYYY@{}}
\toprule
统计/门槛 & Keywords/数量 & KL/条件 & 说明 \\
\midrule
最小值 & 0.54 & 0.000905 & 全部 120 次 \\
线性 25 分位 & 0.88 & 0.003696 & 全部 120 次 \\
中位数 & 0.98 & 0.010354 & 全部 120 次 \\
线性 75 分位 & 1.00 & 0.023271 & 全部 120 次 \\
最大值 & 1.00 & 0.115156 & 全部 120 次 \\
Trial 66 & 0.54 & 0.089975 & 描述性示例；未通过验收 \\
Trial 119 & 0.56 & 0.031919 & 描述性示例；未通过验收 \\
COMPLETE & 120/120 & 非 COMPLETE：0 & 单次自适应搜索 \\
Keywords 上限 & 0 & K <= 0.10 & 通过数 \\
相对下降 & 0 & (K0-K)/K0 >= 0.50 & 通过数 \\
KL 上限 & 120 & D <= 0.15 & 通过数 \\
联合验收 & 0 & failed & 未选中候选；审计等结果为空 \\
Pareto 点 & 10 & 双目标均最小化 & 相同坐标不互相支配 \\
\bottomrule
\end{tabularx}
\end{table*}


表~\ref{tab:ara-v1-results} 给出全部候选的线性分位数。Keywords 的最小值、中位数和最大值分别为 0.54、0.98、1.00，KL 对应的边际统计为 0.000905、0.010354、0.115156。不同指标的边际极值通常来自不同候选，不能组合成一个不存在的最佳模型。所有 KL 值低于门槛仅说明首词元分布代理满足设定上限。

最低 Keywords 的描述性例子为 ID 66：$K=0.54$、$D=0.089975$，解析层区间 $[16,52)$，优化 72 个模块。其相对关键词下降为 0.46；因本次基线恰为 1.00，也等于下降 46 个百分点，两种量在一般情况下不可互换。该候选既未达到 0.10 的绝对门槛，也未达到 0.50 的相对下降门槛，不是被验收或导出的模型。末个试次 ID 119 为 $K=0.56$、$D=0.031919$；终端从 1 开始计数，显示的第 120 次试验对应从 0 开始计数的 journal ID 119。

\subsection{Pareto 集与搜索轨迹}

在同时最小化 $(D,K)$ 时，若一个候选两项均不差且至少一项严格更好，它才支配另一个候选；重复坐标不互相支配。由全部 120 个点得到十个 Pareto ID：28、40、66、79、94、100、101、112、118、119。非支配性仅相对于本次有限候选集成立，不等于通过式~\eqref{eq:v1-gate}，也不是总体 Pareto 最优性证书。

\begin{figure*}[t]
\centering
\begin{minipage}[t]{0.49\textwidth}
\centering
\begin{tikzpicture}
\begin{axis}[
 width=\linewidth,height=57mm,
 xlabel={首词元 KL（验证集）},ylabel={Keywords},
 xmin=0,xmax=0.16,ymin=0,ymax=1.05,
 xtick={0,0.05,0.10,0.15},ytick={0,0.1,0.5,1},
 scaled ticks=false,tick label style={font=\scriptsize},
 label style={font=\small},grid=major,grid style={black!8},
 legend style={font=\scriptsize,draw=none,at={(0.5,1.02)},
 anchor=south,legend columns=2},
 scatter/classes={0={mark=*,draw=black!45,fill=black!25},
 1={mark=diamond*,draw=red!70!black,fill=red!60}},
]
\addplot[scatter,only marks,scatter src=explicit symbolic,mark size=1.5pt,forget plot]
 table[x=kl_divergence,y=keywords,meta=pareto,col sep=comma]
 {evidence/ara-v1-trials.csv};
\addlegendimage{only marks,mark=*,black!45}
\addlegendentry{其余候选}
\addlegendimage{only marks,mark=diamond*,red!70!black}
\addlegendentry{Pareto 候选}
\addplot[black,dashed,forget plot] coordinates {(0.15,0) (0.15,1.05)};
\addplot[black,dashed,forget plot] coordinates {(0,0.10) (0.16,0.10)};
\node[anchor=south west,font=\scriptsize] at (axis cs:0.005,0.11)
 {$K\leq0.10$};
\end{axis}
\end{tikzpicture}
\par\small (a) 所有候选与 Pareto 集
\end{minipage}\hfill
\begin{minipage}[t]{0.49\textwidth}
\centering
\begin{tikzpicture}
\begin{axis}[
 width=\linewidth,height=57mm,
 xlabel={试次 ID（从 0 开始）},ylabel={Keywords},
 xmin=0,xmax=119,ymin=0,ymax=1.05,
 xtick={0,35,66,90,119},ytick={0,0.1,0.5,1},
 tick label style={font=\scriptsize},label style={font=\small},
 grid=major,grid style={black!8},
 legend style={font=\scriptsize,draw=none,at={(0.5,1.02)},
 anchor=south,legend columns=2},
]
\addplot[only marks,mark=*,mark size=1.2pt,black!45]
 table[x=trial_number,y=keywords,col sep=comma]
 {evidence/ara-v1-trials.csv};
\addlegendentry{每次 Keywords}
\addplot[blue!70!black,thick,const plot]
 table[x=trial_number,y=running_best_keywords,col sep=comma]
 {evidence/ara-v1-trials.csv};
\addlegendentry{累计最低值}
\addplot[black,dashed,forget plot] coordinates {(35.5,0) (35.5,1.05)};
\addplot[black,dotted,forget plot] coordinates {(0,0.10) (119,0.10)};
\node[anchor=north,font=\scriptsize] at (axis cs:17,0.44) {startup};
\node[anchor=north,font=\scriptsize] at (axis cs:78,0.44) {自适应 TPE};
\end{axis}
\end{tikzpicture}
\par\small (b) 搜索顺序与累计最低值
\end{minipage}
\caption{point-v1 在验证集上的一次自适应搜索，最终验收失败。
(a) 全部 120 个候选；红色标出十个非支配点，虚线为 KL 0.15 和 Keywords 0.10 门槛。
(b) 同一验证集的逐次关键词率与累计最低值，竖线位于 ID 35 之后，区分 36 次 startup 与其余 TPE 试次。
图中不含独立重复、平滑估计或误差区间；Pareto 候选也未通过联合门槛。}
\label{fig:ara-v1-search}
\end{figure*}


图~\ref{fig:ara-v1-search} 同时呈现二维代理空间与按试次顺序的累计最低 Keywords。startup 后出现更低关键词候选，但这是同一次自适应搜索的描述，不能据此推断 TPE 相对其他策略的因果优势。累计最小值必然单调不增，也不是每次更新均改善的证据。试次 66 之后未刷新 0.54 的最低记录，不证明扩大预算必然无效，更不证明某个新目标必然有效。

\subsection{执行观察与复现限制}

日志报告总耗时 2 小时 38 分钟，终端采样的 GPU 已分配显存约 52.05 GB、常驻系统内存约 5.47 GB。这里保留日志单位及 allocated/resident 含义，不将其转换为完整运行峰值或保留显存。120 次完成支持本次未记录数值故障的有限观察；没有修复前对照、跨种子或跨模型重复，不能估计一般 OOM 概率或归因规范化修复的稳定性收益。

进程退出码为零，外层脚本仍打印成功及适配器路径，但机器验收明确失败。适配器未导出依据归档运行报告，不能只由本地目录缺失推断远端状态。本文交付的图表与 PDF 不构成模型适配器的成功导出证据。

运行记录的 HEAD 为 \code{cd2977a}，归档报告说明当时部署了未提交修复，后续对应修复提交为 \code{868ca73}。后者并非精确运行工作树的完整快照，仅检出前者也无法重建已部署修复。journal 的 \code{model\_commit} 为 \code{null}，模型指纹不能替代可恢复的模型快照，不得倒填 v2 模板的模型修订；未记录的包版本继续列为缺失。

归档没有逐提示响应、logits 或逐样本 KL 数组，因而不能重新判定词汇命中、语义拒答或散度。它只允许对已存评分的聚合进行一致性检查。单一模型、单次搜索与反复使用的验证集限制了外推；没有独立审计结果，不能报告跨种子置信区间、语义攻击成功率、全面能力保持或机制因果性。
